\documentclass{article}
\usepackage{mathnotes}

\title{Conway's Game of Life}
\description{Conway's famous cellular automaton demonstrating emergence and complexity from simple rules, featuring still lifes, oscillators, spaceships, and computational universality with interactive simulation.}

\begin{document}

\section{Conway's Game of Life}

Conway's Game of Life is a cellular automaton devised by mathematician John Horton Conway in 1970. Despite its simple rules, it exhibits remarkably complex behavior and has become one of the most well-known examples of emergence in computational systems.

\subsection{Rules}

The Game of Life is played on a two-dimensional grid where each cell can be either alive (1) or dead (0). The state of each cell evolves according to four simple rules:

\begin{enumerate}
\item \textbf{Underpopulation}: Any live cell with fewer than two live neighbors dies
\item \textbf{Survival}: Any live cell with two or three live neighbors survives
\item \textbf{Overpopulation}: Any live cell with more than three live neighbors dies
\item \textbf{Reproduction}: Any dead cell with exactly three live neighbors becomes alive
\end{enumerate}

These rules are applied simultaneously to all cells in each generation.

\subsection{Interactive Simulation}

Explore the Game of Life with this interactive demonstration. You can:

\begin{itemize}
\item Click cells to toggle them alive/dead
\item Use the controls to start, stop, and step through generations
\item Adjust the simulation speed and initial population
\item Watch population statistics in real-time
\end{itemize}

\includedemo{game-of-life}

\subsection{Common Patterns}

\subsubsection{Still Lifes}
Patterns that don't change:

\begin{itemize}
\item \textbf{Block}: 2×2 square
\item \textbf{Beehive}: 6 cells in a hexagonal shape
\item \textbf{Loaf}: 7 cells in a distinctive shape
\item \textbf{Boat}: 5 cells resembling a boat
\end{itemize}

\subsubsection{Oscillators}
Patterns that cycle through a fixed sequence:

\begin{itemize}
\item \textbf{Blinker}: Period 2, alternates between 3 horizontal and 3 vertical cells
\item \textbf{Toad}: Period 2, 6 cells that shift shape
\item \textbf{Beacon}: Period 2, two blocks connected diagonally
\item \textbf{Pulsar}: Period 3, a spectacular 48-cell pattern
\end{itemize}

\subsubsection{Spaceships}
Patterns that translate across the grid:

\begin{itemize}
\item \textbf{Glider}: 5 cells that move diagonally
\item \textbf{Lightweight Spaceship (LWSS)}: Moves horizontally
\item \textbf{Middleweight Spaceship (MWSS)}: Larger horizontal mover
\item \textbf{Heavyweight Spaceship (HWSS)}: Even larger horizontal mover
\end{itemize}

\subsection{Mathematical Properties}

\subsubsection{Computational Universality}
The Game of Life is Turing complete, meaning it can simulate any computation that can be described algorithmically. This was proven by constructing:

\begin{itemize}
\item Logic gates (AND, OR, NOT)
\item Memory storage
\item Signal transmission
\end{itemize}

\subsubsection{Growth Patterns}
Starting from finite patterns, populations can:

\begin{enumerate}
\item Die out completely
\item Stabilize at a constant population
\item Oscillate between fixed states
\item Grow indefinitely (rare but possible)
\end{enumerate}

\subsubsection{Density Classification}
The critical density for random initial configurations is approximately 0.37. Above this, patterns tend to die out; below it, they tend to stabilize at lower densities.

\subsection{Notable Discoveries}

\subsubsection{Gosper Glider Gun}
Discovered by Bill Gosper in 1970, this was the first pattern found that grows indefinitely, emitting a stream of gliders.

\subsubsection{Garden of Eden}
Patterns that cannot arise from any previous configuration through the rules of the game.

\subsubsection{Methuselahs}
Small patterns that take a long time to stabilize:

\begin{itemize}
\item \textbf{R-pentomino}: 5 cells that evolve for 1103 generations
\item \textbf{Acorn}: 7 cells that stabilize after 5206 generations
\item \textbf{Diehard}: 7 cells that vanish after 130 generations
\end{itemize}

\subsection{Applications and Influence}

The Game of Life has influenced many fields:

\begin{enumerate}
\item \textbf{Computer Science}: Demonstrating emergence and self-organization
\item \textbf{Biology}: Modeling population dynamics and pattern formation
\item \textbf{Physics}: Studying phase transitions and critical phenomena
\item \textbf{Philosophy}: Exploring questions about determinism and complexity
\item \textbf{Art}: Creating generative and algorithmic art
\end{enumerate}

\subsection{Implementation Notes}

Efficient implementations often use:

\begin{itemize}
\item Sparse data structures for large, mostly empty grids
\item Bit manipulation for parallel cell updates
\item HashLife algorithm for computing far-future states
\item GPU acceleration for real-time visualization
\end{itemize}

\subsection{Further Exploration}

\begin{itemize}
\item \href{http://www.conwaylife.com/wiki/}{LifeWiki}: Comprehensive database of Life patterns
\item Gardner, M. (1970). "The fantastic combinations of John Conway's new solitaire game 'life'"
\item Berlekamp, E., Conway, J., Guy, R. (2001). \emph{Winning Ways for Your Mathematical Plays}
\end{itemize}

\end{document}
